% Bibliografía del capítulo de confiabilidad (verificada en línea el 2 de octubre de 2026)
\bibitem{confiabilidad-vesely1981} W.~E.~Vesely, F.~F.~Goldberg, N.~H.~Roberts y
  D.~F.~Haasl. \emph{Fault Tree Handbook}. NUREG-0492, U.S. Nuclear Regulatory
  Commission, Washington, 1981.

\bibitem{confiabilidad-stamatelatos2002} M.~Stamatelatos, W.~Vesely, J.~Dugan,
  J.~Fragola, J.~Minarick y J.~Railsback. \emph{Fault Tree Handbook with Aerospace
  Applications}, versión~1.1. NASA Office of Safety and Mission Assurance, Washington,
  2002.

\bibitem{confiabilidad-iec60812} IEC 60812:2018. \emph{Failure modes and effects
  analysis (FMEA and FMECA)}, 3.ª ed. International Electrotechnical Commission, Ginebra,
  2018.

\bibitem{confiabilidad-iec61078} IEC 61078:2016. \emph{Reliability block diagrams},
  3.ª ed. International Electrotechnical Commission, Ginebra, 2016.

\bibitem{confiabilidad-rausand2004} M.~Rausand y A.~Høyland. \emph{System Reliability
  Theory: Models, Statistical Methods, and Applications}, 2.ª ed. Wiley, Hoboken, 2004.

\bibitem{confiabilidad-aiagvda2019} AIAG y VDA. \emph{FMEA Handbook}, 1.ª ed. Automotive
  Industry Action Group y Verband der Automobilindustrie, 2019.

\bibitem{confiabilidad-milstd882e} U.S. Department of Defense. \emph{MIL-STD-882E:
  Standard Practice for System Safety}, 11 de mayo de 2012.

\bibitem{confiabilidad-faa2021} Federal Aviation Administration. Operation of Small
  Unmanned Aircraft Systems Over People. \emph{Federal Register}, documento 2020-28947,
  15 de enero de 2021. Requisitos de la categoría~2 en 14 CFR \S\,107.120.

\bibitem{confiabilidad-feng2018} X.~Feng, M.~Ouyang, X.~Liu, L.~Lu, Y.~Xia y X.~He.
  Thermal runaway mechanism of lithium ion battery for electric vehicles: A review.
  \emph{Energy Storage Materials}, 10:246--267, 2018. doi:10.1016/j.ensm.2017.05.013.
